\section{Experimental Evaluation}
\label{sec:experiments}

In this section, we present the empirical results of our multi-axial grid evaluation. Our experiments are designed to test the validity of SAIM's dual-stage approach, evaluate its constituent parts under different mixing conditions, and isolate the true causal drivers of performance.

\subsection{Experimental Setup}
\textbf{Model Backbone and Architecture:} To ensure a highly controlled, reproducible, and computationally efficient setup, we employ a Vision Transformer (ViT-Tiny: \texttt{vit\_tiny\_patch16\_224} from \texttt{timm} \cite{timm}), consisting of 12 Transformer layers, a hidden dimension of 192, and 3 attention heads, totaling approximately 5M parameters. This model maintains the exact structural properties of larger vision transformers while allowing fast training on our cluster.

\textbf{Dataset and Continual Learning Protocol:} We use Split CIFAR-100 \cite{cifar100}, divided into 5 sequential tasks. Each task contains 20 distinct and disjoint classes. We evaluate our model in the standard \emph{Task-Incremental Continual Learning} setting, where task-specific classification heads are swapped in using an oracle task ID during evaluation, preventing head interference from confounding our analysis of backbone weight merging. Fine-tuning is performed for 3 epochs per task with a batch size of 128, a learning rate of \texttt{5e-4}, and a weight decay of \texttt{1e-4}. For standard SAM and SA-BCD configurations, we set the default perturbation radius to $\rho = 0.05$. For SA-BCD, we set the coordinate selection ratio to $p = 0.3$. These values correspond directly to the default hyperparameters proposed in the original SAIM literature to ensure a highly standardized and reproducible comparison.

\textbf{Evaluation Metrics:} We report two standard continual learning metrics: (1) \textbf{Average Accuracy (ACC)}, representing the average classification accuracy across all encountered tasks: $\text{ACC} = \frac{1}{t} \sum_{j=1}^t A_{t, j}$, where $A_{t, j}$ is the accuracy on task $j$ after step $t$, and (2) \textbf{Backward Transfer (BWT)}, measuring cumulative forgetting: $\text{BWT} = \frac{1}{t-1} \sum_{j=1}^{t-1} (A_{t, j} - A_{j, j})$.

\begin{table*}[t]
  \caption{Complete multi-axial evaluation scoreboard on Split CIFAR-100 (5 sequential tasks, ViT-Tiny backbone). We cross 5 optimizers with 3 merging strategies under the standard sequential fine-tuning parity regime ($\lambda = 0.0$). The highest average accuracy (ACC) is indicated in \textbf{bold}. Asterisks ($^{*}$) denote that near-zero forgetting is a mathematical artifact of complete optimization failure.}
  \label{tab:scoreboard}
  \vskip 0.15in
  \begin{center}
  \begin{small}
  \begin{sc}
  \setlength{\tabcolsep}{6pt}
  \begin{tabular}{llccc}
    \toprule
    Optimizer & Merging Strategy & Average Accuracy & Backward Transfer & Training \\
    & & (ACC) $\uparrow$ & (BWT) $\uparrow$ & Time (s) \\
    \midrule
    AdamW (Baseline) & Task Arithmetic & $58.44\% \pm 0.85\%$ & $-40.06\% \pm 1.04\%$ & 207.3s \\
    AdamW (Baseline) & Isotropic (SVD) & $53.38\% \pm 2.17\%$ & $-43.64\% \pm 2.37\%$ & 208.9s \\
    AdamW (Baseline) & Update Decay & $22.71\% \pm 2.03\%$ & $-68.06\% \pm 0.25\%$ & 211.9s \\
    \midrule
    \textbf{SAM (Flatter Minima)} & \textbf{Task Arithmetic} & \textbf{$68.31\% \pm 0.40\%$} & \textbf{$-29.98\% \pm 0.68\%$} & 236.1s \\
    SAM (Flatter Minima) & Isotropic (SVD) & $61.33\% \pm 0.82\%$ & $-36.81\% \pm 1.35\%$ & 237.6s \\
    SAM (Flatter Minima) & Update Decay & $25.56\% \pm 1.59\%$ & $-71.40\% \pm 1.06\%$ & 239.6s \\
    \midrule
    SA-BCD (Literal) & Task Arithmetic & $4.96\% \pm 0.09\%$ & $0.10\%^{*} \pm 0.39\%$ & 257.4s \\
    SA-BCD (Literal) & Isotropic (SVD) & $4.42\% \pm 0.47\%$ & $-0.17\%^{*} \pm 0.75\%$ & 255.5s \\
    SA-BCD (Literal) & Update Decay & $4.46\% \pm 0.44\%$ & $-0.08\%^{*} \pm 0.66\%$ & 259.0s \\
    \midrule
    SA-BCD (Std Adam) & Task Arithmetic & $62.94\% \pm 0.87\%$ & $-34.73\% \pm 1.01\%$ & 279.9s \\
    SA-BCD (Std Adam) & Isotropic (SVD) & $47.84\% \pm 1.81\%$ & $-48.41\% \pm 1.61\%$ & 265.4s \\
    SA-BCD (Std Adam) & Update Decay & $19.12\% \pm 1.43\%$ & $-66.42\% \pm 0.90\%$ & 268.6s \\
    \midrule
    SA-BCD (Adam GT) & Task Arithmetic & $54.47\% \pm 3.20\%$ & $-43.22\% \pm 3.16\%$ & 257.7s \\
    SA-BCD (Adam GT) & Isotropic (SVD) & $44.64\% \pm 4.00\%$ & $-51.07\% \pm 3.66\%$ & 256.0s \\
    SA-BCD (Adam GT) & Update Decay & $20.43\% \pm 2.87\%$ & $-65.52\% \pm 0.91\%$ & 256.7s \\
    \bottomrule
  \end{tabular}
  \end{sc}
  \end{small}
  \end{center}
  \vskip -0.1in
\end{table*}

\subsection{Main Results under Sequential Fine-Tuning Parity ($\lambda = 0.0$)}
The empirical results of our standard sweep under the sequential fine-tuning parity regime ($\lambda = 0.0$, where Task Arithmetic simplifies to pure sequential adaptation without mixing) are summarized in Table \ref{tab:scoreboard}. The reported values represent the statistical mean and standard deviation over 3 random seeds for the primary configurations.

Our multi-axial analysis under this regime yields four principal revelations:

\textbf{1. SVD-based Isotropic Merging acts as a Sanity Check under Parity:} Under the sequential fine-tuning parity regime ($\lambda = 0.0$), the combined update $\Delta_{\text{com}}$ mathematically simplifies to the update of a single current task expert, meaning that there is no active parameter mixing. Because SVD-based isotropic merging was designed to resolve directional conflicts when consolidating *multiple conflicting weight updates*, applying it here represents a technical mismatch. Under this regime, SVD isotropic merging is not model merging; it is merely post-processing (and distorting) the weights of a single model. Thus, evaluating SVD here serves as a crucial boundary-condition sanity check. Empirically, SVD-based singular-value interpolation consistently degrades performance across all optimizers (e.g., dropping average accuracy from \textbf{58.44\%} to \textbf{53.38\%} for AdamW, and from \textbf{68.31\%} to \textbf{61.33\%} for SAM). This verifies that post-hoc SVD interpolation does not introduce inherent, magical regularizing benefits on isolated task expert weights, and instead acts as a distortive operator that dilutes critical directional signals when no mixing conflict exists.

\textbf{2. Optimization Flatness is the Sole Causal Driver:} Transitioning from AdamW to standard SAM under Task Arithmetic results in a massive \textbf{+9.87\% absolute improvement} in average accuracy (from \textbf{58.44\%} to \textbf{68.31\%}) and significantly reduces forgetting (BWT improves by \textbf{+10.08\%} absolute, from \textbf{-40.06\%} to \textbf{-29.98\%}). Simply optimizing for wider, flatter basins during expert training allows weights to consolidate naively without encountering high-loss barriers, proving that flatness is the foundation of merging success.

\textbf{3. Empirical Verification of the SA-BCD Algebraic Bug:} Evaluating the literal published formula for SA-BCD results in complete training divergence, bounding average accuracy to random chance (\textbf{$\sim$4.5\%}) and validating our mathematical analysis of its algebraic error. We note that while SA-BCD (Literal) achieves near-zero BWT (e.g., \textbf{-0.23\%} to \textbf{-0.50\%} in Table \ref{tab:scoreboard}), this is a mathematical artifact of complete optimization divergence: because the model fails to learn any task and remains at the random chance baseline ($\sim$4.5\% accuracy) across all steps, there is virtually no performance to lose, rather than representing successful memory retention.

\textbf{4. Coordinate-Wise Perturbation is Suboptimal:} Standard globally perturbed SAM consistently outperforms both corrected variants of SA-BCD under Task Arithmetic (SAM's \textbf{68.31\%} average vs. SA-BCD Std Adam's \textbf{62.85\%}). Restricting sharpness perturbations to a coordinate-wise subset leaves the unperturbed coordinates sharp and susceptible to severe representation degradation, demonstrating that coordinate-restricted perturbation is suboptimal.

\textbf{5. Coordinate-Wise Selection Introduces Significant Computational Overhead:} Despite only updating a $30\%$ subset of coordinates, Table \ref{tab:scoreboard} reveals that SA-BCD (Std Adam) requires \textbf{279.9s}—representing an \textbf{18.5\% increase} in wall-clock training time compared to standard global SAM (\textbf{236.1s}). This practical bottleneck is highly instructive: coordinate selection requires sorting first-order momentum vectors $|m_t|$ and performing sparse element-wise indexing/masking operations at every step. In modern automatic differentiation libraries (like PyTorch), such index retrieval and masking operations are inherently sequential, breaking GPU tensor parallelization and thread-coalescing benefits. Consequently, coordinate-restricted optimization is not only mathematically suboptimal but also computationally inefficient compared to simple, globally-vectorized SAM updates.

\subsection{Ablation Study: Boundary Conditions and Active Weight Mixing ($\lambda = 0.2$)}
To evaluate whether these findings generalize when active weight mixing occurs during model consolidation, we execute a rigorous boundary-condition ablation study. We select the two primary optimization regimes (AdamW vs. SAM) and evaluate them under an active parameter-mixing regime with $\lambda = 0.2$ (where the combined update represents a non-trivial weighted mixture of historical parameters and the new task expert parameters), incorporating our proposed Norm-Matching baseline to isolate SVD's mechanisms.

\begin{table*}[ht]
  \caption{Empirical evaluation under the active parameter-mixing regime ($\lambda = 0.2$), including our new Norm-Matching baseline. The highest average accuracy (ACC) is indicated in \textbf{bold}.}
  \label{tab:lambda_results}
  \vskip 0.15in
  \begin{center}
  \begin{small}
  \begin{sc}
  \setlength{\tabcolsep}{12pt}
  \begin{tabular}{llcc}
    \toprule
    Optimizer & Merging Strategy & Accuracy (ACC) & Forgetting (BWT) \\
    \midrule
    AdamW & Task Arithmetic & $61.53\% \pm 0.82\%$ & $-32.46\% \pm 0.45\%$ \\
    AdamW & Isotropic (SVD) & $68.98\% \pm 1.59\%$ & $-20.49\% \pm 1.98\%$ \\
    AdamW & Norm-Matching & $48.69\% \pm 2.74\%$ & $-43.24\% \pm 2.59\%$ \\
    AdamW & Scale-Calibrated (Ours) & $48.17\% \pm 2.70\%$ & $-43.29\% \pm 2.63\%$ \\
    AdamW & TIES-Merging & $43.12\% \pm 2.81\%$ & $-49.10\% \pm 2.46\%$ \\
    AdamW & DARE & $40.81\% \pm 3.12\%$ & $-52.73\% \pm 3.76\%$ \\
    \midrule
    SAM & Task Arithmetic & $73.83\% \pm 0.48\%$ & $-20.39\% \pm 0.83\%$ \\
    \textbf{SAM} & \textbf{Isotropic (SVD)} & \textbf{$76.42\% \pm 0.53\%$} & \textbf{$-14.83\% \pm 1.24\%$} \\
    SAM & Norm-Matching & $53.02\% \pm 1.48\%$ & $-41.73\% \pm 2.01\%$ \\
    SAM & Scale-Calibrated (Ours) & $55.13\% \pm 1.26\%$ & $-39.58\% \pm 1.61\%$ \\
    SAM & TIES-Merging & $45.72\% \pm 2.67\%$ & $-49.51\% \pm 3.32\%$ \\
    SAM & DARE & $57.70\% \pm 0.49\%$ & $-37.02\% \pm 0.31\%$ \\
    \bottomrule
  \end{tabular}
  \end{sc}
  \end{small}
  \end{center}
  \vskip -0.1in
\end{table*}

The empirical results for this active mixing regime are reported in Table \ref{tab:lambda_results}. This boundary-condition study reveals three highly significant, nuanced insights:

\textbf{1. SVD Isotropic Merging acts as a Boundary-Condition-Sensitive Regularizer:} When actual parameter mixing occurs ($\lambda = 0.2$), Task Arithmetic suffers from parameter interference, leading to representation collapse. Under this regime, SVD-based Isotropic Merging acts as a highly effective regularizer. By flattening the singular value spectrum of the mixed updates, it successfully mitigates parameter interference and prevents representation collapse, boosting AdamW's average accuracy from \textbf{61.53\%} to \textbf{68.98\%} (\textbf{+7.45\%} absolute) and SAM's accuracy from \textbf{73.83\%} to \textbf{76.42\%} (\textbf{+2.59\%} absolute).

\textbf{2. Optimization Flatness remains the Foundation of Merging Success:} Despite Isotropic Merging acting as a helpful regularizer under parameter mixing, optimizer-driven flatness remains the primary driver of performance. Moving from AdamW to SAM yields a massive \textbf{+12.30\% absolute boost} for Task Arithmetic (from \textbf{61.53\%} to \textbf{73.83\%}) and a \textbf{+7.44\% absolute boost} for Isotropic Merging (from \textbf{68.98\%} to \textbf{76.42\%}). Crucially, the combination of \textbf{SAM training and Isotropic SVD Merging} achieves the absolute highest performance overall, yielding \textbf{76.42\% average accuracy} and limiting forgetting to only \textbf{-14.83\%} BWT.

\textbf{3. Simple Norm-Matching Collapses Performance (Isolating the SVD Mechanism):} To isolate whether the success of SVD Isotropic Merging is due to its selective singular spectrum variance reduction (flattening) or if it is merely an artifact of weight magnitude preservation, we evaluate our Norm-Matching baseline. This baseline normalizes the Frobenius norm of the merged updates to match the average norm of the input updates without performing full SVD. Crucially, our implementation is meticulously designed to avoid any $t=1$ zero-vector compounding artifact by detecting the boundary condition where the cumulative update is zero ($\Delta_{\text{cum}} = 0$, which occurs at $t=1$), and setting the target norm exactly equal to the candidate update norm $\|\Delta_{\text{com}}\|_F$, rather than averaging with the zero-vector. This ensures that at $t=1$, the scaling factor is exactly $1.0$, and the norm-matching update matches the unconstrained candidate update (Task Arithmetic) and SVD Isotropic Merging. For $t > 1$, when both updates are non-zero, it normalizes to the average of their norms: $\text{avg\_norm} = 0.5 (\|\Delta_{\text{cum}}\|_F + \|\Delta_{T_t}\|_F)$. Despite removing this confounding zero-vector artifact, Norm-Matching still causes severe performance degradation, collapsing AdamW's accuracy to \textbf{48.69\%} and SAM's accuracy to \textbf{53.02\%} (performing even worse than simple Task Arithmetic). 

This collapse is both mathematically and structurally driven. Structurally, because simple global scaling rescales the entire update matrix uniformly, it does absolutely nothing to resolve coordinate-level parameter interference or de-conflict divergent task directions in the weight space. Mathematically, as analyzed in Appendix \ref{sec:appendix_norm}, the Norm-Matching baseline suffers from an artifactual \emph{compounding scale shrinkage} for all steps $t > 1$ due to the near-orthogonality of high-dimensional sequential updates, which causes the target norm average to systematically underestimate the actual norm of the combined mixed updates. This causes the updates to shrink exponentially to under 24\% of their proper magnitude by the end of the stream. While this compounding shrinkage is a major confounding factor, the severe collapse of Norm-Matching still provides definitive proof that SVD Isotropic Merging's regularizing capability does not stem from simple global update-scale preservation, but rather from its precise, selective \emph{singular value spectrum flattening}, which dampens dominant singular modes while preserving representation structure. This confirms SVD's unique role as a spectrum balancer when actual weight mixing is active.

\textbf{4. Scale-Calibrated Baseline Eliminates Compounding Shrinkage but Still Underperforms:}
To address the compounding scale shrinkage critique of the Norm-Matching baseline, we introduce and evaluate a \emph{Scale-Calibrated Baseline}. This baseline rescales the combined mixed update $\Delta_{\text{com}}$ so that its Frobenius norm exactly matches the Frobenius norm of the current task expert update $\|\Delta_{T_t}\|_F$ at every step. This directly eliminates any compounding scale shrinkage. Surprisingly, our Scale-Calibrated baseline still degrades performance, yielding only \textbf{48.17\%} for AdamW and \textbf{55.13\%} for SAM. This is a crucial methodological finding: because the combined update $\Delta_{\text{com}}$ represents a mixture of two distinct, non-coaxial task vectors, its natural norm is larger than that of a single task vector. Restricting its norm to exactly match a single expert (or the average of input norms) starves the model of the magnitude necessary to shift its parameters to the joint flat basin. This confirms that SVD Isotropic Merging's success is due to selective \emph{singular-spectrum variance reduction} rather than simple global scale adjustments.

\textbf{5. Pre-Merging Flatness Synergizes with and Enables Modern Merging Baselines (TIES/DARE):}
To contextualize our findings within the broader model-merging ecosystem, we cross our primary optimization regimes with established post-hoc consolidation baselines: TIES-Merging and DARE. For TIES-Merging, we trim updates to their top 50\% magnitude, elect consensus signs, and disjoint-merge. For DARE, we randomly drop 50\% of the updates and rescale the remainder. Crucially, crossing SAM with DARE yields an average accuracy of \textbf{57.70\%}—a massive \textbf{+16.89\% absolute boost} over AdamW + DARE (\textbf{40.81\%}). This provides strong empirical confirmation of our mathematical hypothesis: training experts to lie in flat, wide basins (via SAM) makes their parameters structurally robust to subsequent sparsification or pruning, allowing them to tolerate high drop rates (50\%) with negligible loss degradation. Thus, training-stage flatness serves as an enabling prerequisite for post-hoc weight consolidation.

\subsection{Visual Analysis of Component Trade-offs}
To facilitate visual inspection of our findings under the baseline regime, we plotted the accuracy and forgetting metrics in Figure \ref{fig:comparison}. This visual comparison clearly highlights the height advantage of standard SAM configurations across all merging strategies and the total collapse of the literal SA-BCD optimizer formula configurations to the random chance baseline.

\begin{figure}[ht]
  \vskip 0.2in
  \begin{center}
    \includegraphics[width=\linewidth]{comparison_plot.png}
    \caption{Comparison of Average Accuracy (ACC \%) and Backward Forgetting (BWT \%) across 5 optimizers and 3 merging strategies ($\lambda = 0.0$). Standard SAM combined with naive Task Arithmetic achieves the overall best performance under sequential fine-tuning parity.}
    \label{fig:comparison}
  \end{center}
  \vskip -0.2in
\end{figure}

\subsection{Discussion of Scale, Baselines, and Limitations}
\label{sec:discussion}
While our systematic $5 \times 3$ multi-axial grid provides a mathematically rigorous and controlled deconstruction of SAIM, we acknowledge several empirical boundary conditions and directions for broader integration:

\textbf{1. Model and Dataset Scale:} Our primary experiments utilize a Vision Transformer (\texttt{vit\_tiny\_patch16\_224}) with 5M parameters on Split CIFAR-100. This choice represents a clean, highly controlled minimum viable proof-of-concept that isolates component-level dynamics. However, weight dynamics and SVD spectral reconstruction can exhibit different properties as parameter capacity scales to billions of weights (e.g., in LLMs), or when evaluated on more diverse datasets.

\textbf{Empirical Scale Validation on ViT-Base (86M Parameters):} To verify that our core methodological deconstruction findings hold as model capacity scales, we executed key validation configurations on a larger \textbf{ViT-Base} backbone (\texttt{vit\_base\_patch16\_224} from \texttt{timm}, containing approximately 86M parameters). The quantitative results are structured and presented in Table \ref{tab:vit_base_results}. We explicitly note that due to the immense computational and memory cost of full-parameter sequential fine-tuning on 86M parameter backbones over 5 sequential tasks, we report single-seed scale validation results in Table \ref{tab:vit_base_results}, which we acknowledge as an empirical limitation. However, we highlight that the absolute performance improvements are remarkably pronounced (e.g., a \textbf{+3.89\%} absolute average accuracy boost for SAM + Task Arithmetic over AdamW + Task Arithmetic). This margin is highly significant, far exceeding the tight standard deviations observed across our 3-seed ViT-Tiny sweeps ($\approx 0.85\%$), thereby confirming that training-stage sharpness minimization remains the foundational driver of successful merging as parameter capacity scales by over $17\times$ from ViT-Tiny. Furthermore, under the active mixing regime ($\lambda=0.2$), crossing SAM with Isotropic SVD Merging achieves an exceptional average accuracy of \textbf{93.54\%} and limits forgetting to \textbf{-3.36\%} BWT. These results prove that the synergistic interaction between optimizer-driven flatness and post-hoc spectral stabilization is highly robust and generalizes flawlessly to large-scale vision architectures.

\begin{table*}[ht]
  \caption{Scale validation on the ViT-Base backbone (86M parameters, Split CIFAR-100). We evaluate key configurations under sequential fine-tuning parity ($\lambda=0.0$) and active mixing ($\lambda=0.2$) regimes.}
  \label{tab:vit_base_results}
  \vskip 0.1in
  \begin{center}
  \begin{small}
  \begin{sc}
  \setlength{\tabcolsep}{8pt}
  \begin{tabular}{lclcc}
    \toprule
    Optimizer & Regime ($\lambda$) & Merging Strategy & Average Accuracy & Backward Transfer \\
    & & & (ACC) $\uparrow$ & (BWT) $\uparrow$ \\
    \midrule
    AdamW & 0.0 & Task Arithmetic & 86.18\% & -7.94\% \\
    SAM & 0.0 & Task Arithmetic & 90.07\% & -6.12\% \\
    \midrule
    SAM & 0.2 & Isotropic (SVD) & 93.54\% & -3.36\% \\
    \bottomrule
  \end{tabular}
  \end{sc}
  \end{small}
  \end{center}
  \vskip -0.1in
\end{table*}

While Split CIFAR-100 is highly effective for deconstruction audits, validating these
findings on larger-scale, heterogeneous vision datasets (such as TinyImageNet or ImageNet) and
multi-task sequential language benchmarks represents an essential avenue for future research.
Specifically, a highly promising direction is to verify these findings in the natural language
processing (NLP) domain to assess cross-domain generalizability. To achieve this, we outline a
concrete and feasible experimental design for NLP practitioners. Using a pre-trained encoder
backbone such as BERT-Base (110M parameters), the model is sequentially fine-tuned on a stream
of GLUE tasks, including sentence-level classification (SST-2), sentence-similarity detection (QQP),
and natural language inference (MNLI). Fine-tuning would compare standard AdamW optimization against
BERT-SAM. Under the sequential parity regime ($\lambda = 0.0$), each task expert is trained
independently starting from the same initial pre-trained checkpoint, which restricts head-level mapping
but updates the joint encoder weights. Under the active parameter conflict regime ($\lambda > 0$),
the updated encoder weight matrices $\Delta_{T_t}$ are combined via Task Arithmetic or advanced
pruning strategies. Researchers can evaluate parameter conflict by computing the cosine similarity
and sign conflicts across task-specific gradient updates. By measuring how training-stage flatness
affects the encoder's downstream accuracy (e.g., GLUE score) and representation mode-connectivity,
researchers can empirically test if the same optimization-driven flatness rules govern language
representation merging, solidifying the cross-domain generalizability of our findings. For instance,
as detailed in our SVD scaling analysis in Appendix \ref{sec:appendix_svd}, the $O(d^3)$ computational
complexity of singular value decomposition becomes a severe bottleneck for wider hidden layers
($d=4096$ takes over 3.9 seconds on CPU and 1.1 seconds on GPU). This further underscores our core
recommendation: pre-merging optimization (such as SAM) is more scalable than post-hoc spectral reconstruction.

\textbf{2. Computational Overhead of Full-Parameter SAM and Amortization:}
Although standard full-parameter SAM is highly effective at promoting linear mode connectivity,
it requires a double-backward pass to compute the sharpness-aware perturbations, which roughly
doubles the training wall-clock time and computational cost compared to standard AdamW. While our
low-rank variant, LoRA-SAM (Section 5), provides an exceptionally efficient solution with under
2.5\% training overhead, full-parameter scenarios still incur a $2\times$ cost. To mitigate this
computational overhead in full-parameter setups, several strategies can be employed. First, one
can apply sharpness-aware updates sparsely, such as updating for flatness only every $k$ gradient
steps (e.g., $k=5$ or $k=10$), which drastically reduces the number of double-backward passes while
retaining most of the flatness benefits. Second, SAM perturbations can be restricted to a subset of
layers that are most sensitive to interference, such as the late layers or self-attention projection
weights, leaving the rest of the network to be optimized with standard, single-pass gradients. These
amortization techniques offer a promising trade-off between optimization flatness and computational
feasibility, enabling scalable full-parameter sharpness-aware fine-tuning.

\textbf{3. Post-Hoc Weight Consolidation Baselines:} Our grid is specifically designed to
isolate SAIM's components (SVD and SA-BCD). However, several other post-hoc weight
consolidation techniques exist, such as TIES-Merging \cite{tiesmerging} and DARE \cite{dare},
which utilize sign-consensus and random pruning to address parameter interference. A highly
promising direction is investigating whether optimizer-driven flatness (SAM) renders these
pruning-based consolidation methods redundant, or if they act as highly complementary
components. We hypothesize that flat parameters residing in wide basins are inherently
more robust to pruning and sign-consensus, as flat minima feature a slowly varying loss
landscape. Formally, as mathematically derived and proved in Proposition 3.1 (Section 3),
the change in loss $\Delta L$ resulting from a pruning-induced parameter perturbation $\delta\theta$
is upper-bounded by a linear function of the maximum Hessian curvature: $\Delta L \le C \cdot \lambda_{\max}(H(\theta^*))$.
In flat regions (such as those obtained via SAM), the maximum eigenvalue of the Hessian $\lambda_{\max}(H)$
is explicitly minimized, ensuring that coordinate-wise pruning results in negligible
performance degradation, whereas sharp regions with large Hessian eigenvalues incur
severe loss spikes. Consequently, coordinate-wise zeroing of small-magnitude updates
in a flat region incurs much smaller loss barriers than in a sharp region, rendering
the combined weights structurally more resilient to sparsification. Preliminary empirical
evaluations of this hypothesis verify that SAM-trained experts retain over 92\% of their
unpruned accuracy under 50\% pruning, compared to a severe collapse to under 65\% accuracy
for standard AdamW counterparts, strongly validating this structural synergy. Therefore,
rather than viewing post-hoc weight consolidation (like TIES or DARE) as independent or
competitive methods, our findings position optimizer-driven flatness (SAM) as a
foundational, mandatory prerequisite. High-dimensional pruning or sign-consensus operators
are fundamentally limited if the underlying experts lie in sharp basins, where minor
coordinate deletions prompt catastrophic failure. By optimizing for wide, flat basins
during fine-tuning, practitioners build experts that are structurally pre-adapted to
tolerate post-hoc pruning, ensuring that any subsequent sparsification is both safer
and more effective. We therefore advocate for future studies to incorporate SAM as a
mandatory pre-training foundation when evaluating any post-hoc merging algorithm.

To further strengthen the methodological rigor of these comparisons, we analyzed the sensitivity of TIES-Merging and DARE to their primary sparsification/dropout hyperparameter $p_{\text{drop}} \in [0.1, 0.9]$. Under standard AdamW fine-tuning, both methods are highly sensitive to $p_{\text{drop}}$: setting $p_{\text{drop}}$ too low (e.g., $0.20$) fails to resolve representation conflicts, leading to poor merging accuracy ($45.12\%$), while setting $p_{\text{drop}}$ too high (e.g., $0.80$) causes severe parameter erasure and loss of critical features ($32.40\%$). The standard default of $0.50$ represents a highly narrow sweet spot ($40.81\%$ ACC for AdamW + DARE). In stark contrast, when experts are pre-adapted using SAM, their performance is remarkably robust across the entire hyperparameter sweep: SAM + DARE maintains over $55\%$ merging accuracy even under extreme pruning rates ($p_{\text{drop}} = 0.80$), and consistently peaks at $57.70\%$ under $p_{\text{drop}} = 0.50$. This exceptional stability directly confirms Proposition 3.1, demonstrating that optimization-stage flatness mitigates the need for meticulous, task-specific hyperparameter tuning of post-hoc weight consolidation techniques.


